\FloatBarrier
\section{سه آزمایش برای مقاله: \lr{EEG} کجا چیزی دارد که سینماتیک ندارد؟}
\label{sec:controls}

\subsection{پرسش}

نتیجه‌ی فصل‌های پیش دو بخش داشت. \lr{E} به‌تنهایی \lr{EMG} را خوب پیش‌بینی می‌کند، ولی به \lr{K} تقریباً چیزی اضافه نمی‌کند. این نتیجه سه پرسش تازه پیش می‌آورد، و جواب آن‌ها تعیین می‌کند که چه ادعایی را می‌توان در یک مقاله مطرح کرد:
\begin{enumerate}
\item \textbf{آیا \lr{E} چیزی بیش از زمان‌بندی تکلیف می‌داند؟} اگر مدلی که فقط می‌داند «الان در کدام مرحله‌ی تکلیف هستیم» به همان دقت \lr{E} برسد، دقت \lr{E} دلیل خوبی بر خواندن فرمان حرکتی نیست.
\item \textbf{آیا نقشه‌ی الکترود \lr{E} با نقشه‌های کلاسیک علوم اعصاب می‌خواند؟} دو نقشه‌ی مرجع مستقل از مدل وجود دارد: \lr{ERD} (کاهش توان آلفا و بتا هنگام حرکت) و \lr{CMC} (هم‌دوسی \lr{EEG} و \lr{EMG}). اگر \lr{E} از قشر حرکتی بخواند، نقشه‌اش باید به این دو شبیه باشد. این همان گامی است که فصل \ref{sec:literature} پیشنهاد کرده بود.
\item \textbf{آیا \lr{EEG} در جایی که سینماتیک هنوز چیزی نمی‌بیند کمک می‌کند؟} سینماتیک حرکت را بعد از وقوعش می‌بیند، اما فعالیت قشر حرکتی پیش از حرکت شروع می‌شود. پس اگر مدل مجبور باشد \lr{EMG} را چند صد میلی‌ثانیه جلوتر پیش‌بینی کند، سهم \lr{EEG} باید بیشتر شود، به‌ویژه درست پیش از شروع حرکت.
\end{enumerate}

هر سه آزمایش در دفترچه‌ی \lr{eeg\_controls.ipynb} و بسته‌ی \lr{eeg\_controls/} پیاده شده و برای هر ۱۲ شرکت‌کننده اجرا شده است. جدول‌ها و شکل‌های این فصل را اسکریپت \lr{v2-report/analysis/eeg\_controls\_results.py} از نتایج ذخیره‌شده می‌سازد و هیچ مدلی را اجرا نمی‌کند.

\paragraph{آزمون‌های گروهی.} هر جا گفته می‌شود اثری «در سطح گروه معنادار» است، آزمون روی ۱۲ شرکت‌کننده انجام شده است: یا \lr{sign-flip} (علامت عدد هر فرد به‌تصادف عوض می‌شود و میانگین دوباره حساب می‌شود)، یا آزمون \lr{Wilcoxon} یک‌طرفه. هر دو فقط فرض می‌کنند افراد از هم مستقل‌اند. با ۱۲ نفر کوچک‌ترین $p$ ممکن حدود \lr{۰٫۰۰۰۲} است. در یک اجرای قبلی با چهار نفر این عدد \lr{۰٫۱۲۵} بود و هیچ آزمون گروهی نمی‌توانست معنادار شود؛ اجرای کامل این محدودیت را برداشته است.

\begin{figure}[htbp]
\centering
\resizebox{\textwidth}{!}{% EEG control tests: the framework of the three experiments (A timing baseline, B reference maps, C forecast horizon).
% Every column has six boxes of the same height, so the rows line up.
\begin{latin}
\begin{tikzpicture}[node distance=4mm and 6mm, font=\footnotesize,
                    cbox/.style={box, text width=4.4cm, minimum height=17mm},
                    cres/.style={result, text width=4.4cm, minimum height=17mm}]
\node[io, text width=14.6cm] (data) {WAY-EEG-GAL, 12 participants, the same train / validation / test series as KG-GT v2:\\
  EEG \quad kinematics and forces \quad EMG envelope (target) \quad task events of every lift};

% column A: task-timing baseline
\node[cbox, below=9mm of data, xshift=-5.1cm, fill=black!4] (a0) {\textbf{A. Task-timing baseline}\\ no EEG, no kinematics};
\node[cbox, below=of a0] (a1) {event times of each lift\\ (LED on, hand start, first touch, \dots)};
\node[cbox, below=of a1] (a2) {Gaussian bumps from $-2$ to $+3$ s around each event:\\ cue, coarse, phase, phase + weight};
\node[cbox, below=of a2] (a3) {ridge regression per muscle on $u=\sqrt{y/s_c}$, $\alpha$ on validation};
\node[cbox, below=of a3] (a4) {test series: $r$ and VAF,\\ the same scoring as E and K};
\node[cres, below=of a4] (a5) {Is E more accurate than knowing the task timing?};

% column B: reference maps
\node[cbox, below=9mm of data, fill=eegc!8, draw=eegc!70!black] (b0) {\textbf{B. Reference maps}\\ no model involved};
\node[cbox, below=of b0] (b1) {raw EEG and raw EMG (4000 Hz)\\ around the lift events};
\node[cbox, below=of b1] (b2) {ERD\% alpha and beta (before first touch, before lift-off, hold);\\ CMC 15--30 Hz in the hold};
\node[cbox, below=of b2] (b3) {Spearman $\rho$ over 32 electrodes against the rollout maps of E};
\node[cbox, below=of b3] (b4) {sign-flip over participants;\\ electrode shuffle for group maps};
\node[cres, below=of b4] (b5) {Does E read the electrodes where the cortex is active?};

% column C: forecast horizon
\node[cbox, below=9mm of data, xshift=5.1cm, fill=kinc!8, draw=kinc!70!black] (c0) {\textbf{C. Forecast horizon}\\ causal KG-GT v2};
\node[cbox, below=of c0] (c1) {causal mask, left-padded stem, causal filters: inputs up to $t$};
\node[cbox, below=of c1] (c2) {target EMG at $t+h$, $h$ = 0, 200, 500, 1000 ms (48 trained models)};
\node[cbox, below=of c2] (c3) {$F-K$ and the residual test of EEG against K};
\node[cbox, below=of c3] (c4) {skill per task phase: before movement, reach, load, hold, release};
\node[cres, below=of c4] (c5) {Does EEG help where kinematics cannot see yet?};

% conclusion
\node[result, below=9mm of b5, text width=14.6cm, fill=black!4] (out) {What does scalp EEG know about the muscles beyond kinematics:\\
  \textbf{when} the movement happens, or \textbf{how much} each muscle works?};

\foreach \c in {a,b,c}{
  \draw[arr] (data.south -| \c0.north) -- (\c0);
  \draw[arr] (\c0)--(\c1); \draw[arr] (\c1)--(\c2); \draw[arr] (\c2)--(\c3); \draw[arr] (\c3)--(\c4); \draw[arr] (\c4)--(\c5);
  \draw[arr] (\c5.south) -- (\c5.south |- out.north);}
\end{tikzpicture}
\end{latin}
}
\caption{چارچوب سه آزمون کنترلی. هر سه از همان داده و همان تقسیم سری‌ها شروع می‌کنند، ولی از سه راه مستقل به یک پرسش می‌رسند.}
\label{fig:ctrl-framework}
\end{figure}

\paragraph{شکل \ref{fig:ctrl-framework} چه می‌گوید.} هر ستون یک آزمون است و هر کدام یک توضیح جایگزین برای دقت \lr{E} را می‌سنجد.
\begin{itemize}
\item \textbf{ستون \lr{A}} نه \lr{EEG} می‌بیند و نه سینماتیک؛ فقط زمان رخدادها را دارد. اگر این ستون به دقت \lr{E} برسد، دقت \lr{E} را می‌توان با «دانستن مرحله‌ی تکلیف» توضیح داد.
\item \textbf{ستون \lr{B}} اصلاً از مدل استفاده نمی‌کند. نقشه‌های \lr{ERD} و \lr{CMC} مستقیم از سیگنال خام ساخته می‌شوند و فقط در مرحله‌ی آخر با نقشه‌ی \lr{rollout} خروجی \lr{E} مقایسه می‌شوند. اگر این دو شبیه باشند، \lr{E} از جای درستی روی سر می‌خواند؛ اگر نه، احتمال \lr{artifact} بیشتر است.
\item \textbf{ستون \lr{C}} خود مدل را تغییر می‌دهد: مدل فقط گذشته را می‌بیند و باید آینده را پیش‌بینی کند. اگر \lr{EEG} اطلاعاتی داشته باشد که سینماتیک هنوز ندارد، سهمش باید با افق بزرگ‌تر شود.
\end{itemize}
پرسش پایین شکل جمع‌بندی هر سه است: \lr{EEG} می‌داند حرکت \emph{کی} اتفاق می‌افتد، یا هر عضله \emph{چقدر} کار می‌کند؟

% ---------------------------------------------------------------------------------------------
\subsection{آزمایش ۱: خط پایه‌ی زمان‌بندی تکلیف}
\label{sec:controls-phase}

\paragraph{ایده.} یک مدل خطی ساده می‌سازیم که هیچ سیگنالی از بدن یا مغز نمی‌بیند و فقط زمان رخدادهای تکلیف را می‌داند. اگر این مدل از \lr{E} بهتر باشد، هر چیزی که \lr{E} درباره‌ی \lr{EMG} می‌داند را می‌توان با زمان‌بندی تکلیف توضیح داد.

\paragraph{ویژگی‌ها.} هشت رخداد هر بلند کردن از داده‌ی اصلی خوانده می‌شود: روشن شدن \lr{LED}، شروع حرکت دست، اولین تماس، تماس هر دو انگشت، \lr{lift-off}، گذاشتن جسم سر جایش، رها کردن و توقف دست. برای هر رخداد، چند تابع زنگوله‌ای (\lr{Gaussian bump}) با فاصله‌های منظم از ۲ ثانیه پیش تا ۳ ثانیه بعد از آن قرار می‌گیرد. هر ویژگی می‌گوید «الان چقدر به فلان لحظه نسبت به فلان رخداد نزدیکیم». مثلاً ویژگی «\lr{LED} به‌علاوه‌ی یک ثانیه» فقط حدود یک ثانیه بعد از روشن شدن چراغ بزرگ است و بقیه‌ی وقت صفر. \lr{ridge} برای هر ویژگی یک وزن یاد می‌گیرد؛ پس پیش‌بینی آن برای هر عضله یک الگوی زمانی ثابت حول هر رخداد است که در هر بلند کردن تکرار می‌شود. چهار نسخه ساخته شد:
\begin{itemize}
\item \textbf{\lr{cue}} (۲۱ ویژگی): فقط روشن شدن \lr{LED}، با گام و پهنای ربع‌ثانیه. این ساده‌ترین چیزی است که یک مغز، یا یک \lr{artifact} بینایی، درباره‌ی زمان می‌داند.
\item \textbf{\lr{coarse}} (۸۸ ویژگی): هر هشت رخداد، هر کدام ۱۱ زنگوله با گام نیم‌ثانیه و پهنای \lr{۰٫۷۵} ثانیه. این نسخه شکل کلی را می‌گیرد ولی جزئیات سریع را نه.
\item \textbf{\lr{phase}} (۱۷۱ ویژگی): هر هشت رخداد، هر کدام ۲۱ زنگوله با گام و پهنای ربع‌ثانیه، به‌علاوه‌ی سه ویژگی روشن و خاموش برای مرحله‌های بلند کردن: دست در حرکت است (از شروع تا توقف حرکت دست)، انگشت‌ها روی جسم‌اند (از اولین تماس تا رها کردن) و جسم در هواست (از \lr{lift-off} تا گذاشتن جسم).
\item \textbf{\lr{phase + weight}} (۵۱۳ ویژگی): همان \lr{phase}، با یک نسخه‌ی جدا برای هر وزن جسم (سبک، متوسط و سنگین). در هر بلند کردن فقط نسخه‌ی وزن همان جسم روشن است و دو نسخه‌ی دیگر صفرند، تا الگوی هر وزن جدا یاد گرفته شود. \lr{E} هیچ‌وقت وزن را نمی‌گیرد، پس این نسخه فقط یک مرجع بالایی است.
\end{itemize}

\fig{ctrl_features_bumps.png}{0.95}{ویژگی‌های خط پایه‌ی زمان‌بندی برای یک بلند کردن واقعی از \lr{P1}. ردیف اول: پوش \lr{EMG} دو عضله و هشت رخداد همان بلند کردن (خط‌چین‌ها رخدادهای بلند کردن‌های قبل و بعدند). ردیف دوم: ۲۱ زنگوله‌ی \lr{cue} حول روشن شدن \lr{LED}؛ چهار تا پررنگ شده‌اند. ردیف سوم و چهارم: زنگوله‌های اولین تماس در \lr{coarse} و \lr{phase}. ردیف آخر: سه ویژگی مرحله که فقط \lr{phase} دارد.}

\paragraph{شکل \ref{fig:ctrl_features_bumps.png} چه می‌گوید.} این شکل نشان می‌دهد یک «ویژگی زمان‌بندی» دقیقاً چیست.
\begin{itemize}
\item \textbf{ردیف اول} یک بلند کردن را با رخدادهایش نشان می‌دهد. فعالیت \lr{AD} کمی بعد از شروع حرکت دست بالا می‌رود و \lr{FDI} درست هنگام تماس قله‌ی کوتاهی دارد. این قله‌ها همیشه در فاصله‌ی تقریباً ثابتی از رخدادها می‌آیند، و همین است که مدل زمان‌بندی از آن استفاده می‌کند.
\item \textbf{ردیف دوم} ۲۱ زنگوله‌ی \lr{cue} را نشان می‌دهد. هر زنگوله در یک فاصله‌ی مشخص از روشن شدن چراغ قله دارد: از ۲ ثانیه پیش تا ۳ ثانیه بعد. چهار زنگوله‌ی پررنگ در فاصله‌ی \lr{−۱}، صفر، \lr{+۱} و \lr{+۲} ثانیه‌اند. زنگوله‌ها با هم کل این بازه‌ی ۵ ثانیه‌ای را می‌پوشانند. پس \lr{ridge} می‌تواند هر الگوی زمانی هموار در این بازه را بسازد، به شرطی که آن الگو در همه‌ی بلند کردن‌ها نسبت به چراغ یکسان باشد.
\item \textbf{ردیف سوم و چهارم} تفاوت \lr{coarse} و \lr{phase} را نشان می‌دهند. زنگوله‌های \lr{coarse} پهن و کم‌تعدادند و فقط تغییرهای آهسته را می‌سازند. زنگوله‌های \lr{phase} باریک‌تر و دو برابر پرتعدادترند و می‌توانند قله‌ای مثل قله‌ی \lr{FDI} هنگام تماس را دنبال کنند.
\item \textbf{ردیف آخر} سه ویژگی مرحله‌ی \lr{phase} است. زنگوله‌ها از ۳ ثانیه بعد از هر رخداد جلوتر نمی‌روند. این سه ویژگی برای نگه‌داشتن‌های طولانی یک سطح ثابت فراهم می‌کنند، هر قدر هم که نگه‌داشتن طول بکشد.
\end{itemize}

\fig{ctrl_features_matrix.png}{0.95}{همه‌ی ویژگی‌های هر نسخه برای چهار بلند کردن پشت سر هم از \lr{P1}. هر ردیف یک ویژگی و هر ستون یک لحظه است؛ رنگ تیره یعنی مقدار بزرگ. ردیف بالا پوش \lr{EMG} \lr{AD} و وزن جسم هر بلند کردن را نشان می‌دهد.}

\paragraph{شکل \ref{fig:ctrl_features_matrix.png} چه می‌گوید.} این شکل ماتریس ورودی کامل \lr{ridge} را برای هر نسخه نشان می‌دهد.
\begin{itemize}
\item \textbf{هر نوار مورب} یک دسته زنگوله است: زنگوله‌های پشت سر هم یکی بعد از دیگری روشن می‌شوند، و چون هر کدام کمی دیرتر قله دارد، در کنار هم یک خط مورب می‌سازند. هر بلند کردن یک خط مورب برای هر رخداد دارد.
\item \textbf{\lr{cue}} فقط یک دسته دارد. \lr{coarse} و \lr{phase} برای هر هشت رخداد یک دسته دارند، که به ترتیب زمانی بلند کردن روی هم چیده شده‌اند. در \lr{phase} سه ردیف مرحله بالای بقیه است.
\item \textbf{\lr{phase + weight}} سه کپی از ماتریس \lr{phase} دارد. در هر بلند کردن فقط کپیِ وزن همان جسم روشن است: دو بلند کردن اول با جسم متوسط در کپی وسط، سومی با جسم سبک در کپی پایین، و چهارمی با جسم سنگین در کپی بالا. پس این مدل می‌تواند برای هر وزن الگوی جدا یاد بگیرد، مثلاً فعالیت بیشتر در نگه‌داشتن جسم سنگین.
\end{itemize}

\paragraph{برازش و سنجش.} یک \lr{ridge regression} برای هر عضله هدف $u=\sqrt{y/s_c}$ را پیش‌بینی می‌کند، یعنی همان هدفی که \lr{v2} می‌بیند. $\alpha$ روی سری‌های \lr{validation} انتخاب می‌شود و سنجش روی همان سری‌های آزمون \lr{v2} و با همان تابع (\lr{standard\_metrics} در نرخ ۵۰۰ هرتز) است. پس اعداد مستقیم با جدول \ref{tab:v2-subj} مقایسه‌پذیرند.

\paragraph{این مدل پیش‌بینی‌کننده‌ی واقعی نیست.} زمان رخدادها از کل ضبط خوانده می‌شود، از جمله رخدادهایی که هنوز اتفاق نیفتاده‌اند. پس این مدل سقف چیزی است که «فقط زمان‌بندی» می‌تواند بگوید، نه رقیب منصفی برای \lr{E}. هدفش جواب دادن به یک پرسش است: آیا دقت \lr{E} به چیزی بیش از این سقف نیاز دارد؟

\begin{table}[htbp]
\centering
\footnotesize
\caption{مدل‌های زمان‌بندی تکلیف در برابر \lr{E} و \lr{K} روی سری‌های آزمون. شش ستون اول $r$ و شش ستون بعد \lr{VAF} (درصد)، هر دو میانگین روی پنج عضله. \lr{cue}: فقط روشن شدن \lr{LED}؛ \lr{coarse}: هشت رخداد با گام نیم‌ثانیه؛ \lr{phase}: هشت رخداد با گام ریز و مرحله‌های تکلیف؛ \lr{+weight}: همان به‌علاوه‌ی وزن جسم. \lr{E} و \lr{K} از مدل غیرعلّی فصل‌های پیش‌اند.}
\label{tab:ctrl-phase}
\resizebox{\textwidth}{!}{%
\begin{tabular}{lcccccccccccc}
\toprule
 & \multicolumn{6}{c}{$r$} & \multicolumn{6}{c}{\lr{VAF} (درصد)} \\ \cmidrule(lr){2-7} \cmidrule(lr){8-13}
شرکت‌کننده & \lr{cue} & \lr{coarse} & \lr{phase} & \lr{+weight} & \lr{E} & \lr{K} & \lr{cue} & \lr{coarse} & \lr{phase} & \lr{+weight} & \lr{E} & \lr{K} \\
\midrule
\lr{P1} & \lr{۰٫۶۷۰} & \lr{۰٫۸۴۴} & \lr{۰٫۸۶۵} & \lr{۰٫۸۷۵} & \lr{۰٫۵۹۷} & \lr{۰٫۹۱۴} & \lr{۴۵٫۱} & \lr{۷۰٫۵} & \lr{۷۳٫۹} & \lr{۷۵٫۷} & \lr{۳۳٫۹} & \lr{۸۳٫۶} \\
\lr{P10} & \lr{۰٫۵۵۰} & \lr{۰٫۸۰۳} & \lr{۰٫۸۳۲} & \lr{۰٫۸۴۲} & \lr{۰٫۷۲۵} & \lr{۰٫۸۹۸} & \lr{۳۱٫۷} & \lr{۶۶٫۴} & \lr{۷۰٫۵} & \lr{۷۲٫۴} & \lr{۵۳٫۶} & \lr{۸۰٫۲} \\
\lr{P11} & \lr{۰٫۶۲۸} & \lr{۰٫۸۴۹} & \lr{۰٫۸۵۳} & \lr{۰٫۸۴۸} & \lr{۰٫۶۴۹} & \lr{۰٫۸۸۴} & \lr{۳۶٫۹} & \lr{۷۲٫۷} & \lr{۷۳٫۳} & \lr{۷۲٫۶} & \lr{۳۶٫۰} & \lr{۷۸٫۱} \\
\lr{P12} & \lr{۰٫۴۶۶} & \lr{۰٫۷۵۲} & \lr{۰٫۷۶۸} & \lr{۰٫۷۹۴} & \lr{۰٫۶۰۴} & \lr{۰٫۸۴۵} & \lr{۲۲٫۶} & \lr{۵۸٫۷} & \lr{۶۰٫۸} & \lr{۶۵٫۱} & \lr{۳۴٫۴} & \lr{۷۲٫۵} \\
\lr{P2} & \lr{۰٫۵۷۷} & \lr{۰٫۸۳۴} & \lr{۰٫۸۵۱} & \lr{۰٫۸۵۵} & \lr{۰٫۷۶۶} & \lr{۰٫۸۸۸} & \lr{۳۲٫۸} & \lr{۶۹٫۹} & \lr{۷۲٫۷} & \lr{۷۳٫۷} & \lr{۵۸٫۳} & \lr{۷۸٫۸} \\
\lr{P3} & \lr{۰٫۶۳۱} & \lr{۰٫۷۶۲} & \lr{۰٫۷۷۹} & \lr{۰٫۷۸۹} & \lr{۰٫۵۹۲} & \lr{۰٫۸۳۷} & \lr{۴۱٫۳} & \lr{۶۰٫۳} & \lr{۶۲٫۹} & \lr{۶۴٫۶} & \lr{۳۰٫۷} & \lr{۶۹٫۸} \\
\lr{P4} & \lr{۰٫۶۷۶} & \lr{۰٫۸۵۱} & \lr{۰٫۸۶۰} & \lr{۰٫۸۷۰} & \lr{۰٫۷۵۹} & \lr{۰٫۹۰۵} & \lr{۴۳٫۴} & \lr{۷۲٫۹} & \lr{۷۴٫۷} & \lr{۷۶٫۲} & \lr{۵۷٫۷} & \lr{۸۱٫۵} \\
\lr{P5} & \lr{۰٫۴۹۱} & \lr{۰٫۷۲۴} & \lr{۰٫۷۳۷} & \lr{۰٫۷۶۲} & \lr{۰٫۶۰۳} & \lr{۰٫۷۹۵} & \lr{۲۷٫۲} & \lr{۵۷٫۴} & \lr{۵۹٫۱} & \lr{۶۲٫۸} & \lr{۳۴٫۴} & \lr{۶۳٫۸} \\
\lr{P6} & \lr{۰٫۵۶۱} & \lr{۰٫۸۲۵} & \lr{۰٫۸۲۷} & \lr{۰٫۸۳۹} & \lr{۰٫۶۲۶} & \lr{۰٫۸۷۷} & \lr{۲۹٫۳} & \lr{۶۸٫۰} & \lr{۶۸٫۵} & \lr{۷۰٫۲} & \lr{۳۵٫۸} & \lr{۷۶٫۷} \\
\lr{P7} & \lr{۰٫۵۳۴} & \lr{۰٫۸۲۴} & \lr{۰٫۸۳۸} & \lr{۰٫۸۷۰} & \lr{۰٫۷۵۶} & \lr{۰٫۹۱۰} & \lr{۲۸٫۰} & \lr{۶۷٫۷} & \lr{۷۰٫۰} & \lr{۷۵٫۲} & \lr{۵۷٫۷} & \lr{۸۲٫۶} \\
\lr{P8} & \lr{۰٫۵۲۳} & \lr{۰٫۶۹۱} & \lr{۰٫۷۱۴} & \lr{۰٫۷۲۷} & \lr{۰٫۵۸۲} & \lr{۰٫۷۹۲} & \lr{۳۱٫۲} & \lr{۵۲٫۶} & \lr{۵۴٫۹} & \lr{۵۶٫۹} & \lr{۳۲٫۷} & \lr{۵۸٫۴} \\
\lr{P9} & \lr{۰٫۶۰۴} & \lr{۰٫۸۴۹} & \lr{۰٫۸۴۷} & \lr{۰٫۸۵۵} & \lr{۰٫۴۷۸} & \lr{۰٫۹۱۸} & \lr{۳۴٫۷} & \lr{۷۲٫۴} & \lr{۷۲٫۱} & \lr{۷۳٫۳} & \lr{۴٫۵} & \lr{۸۳٫۵} \\
\midrule
\textbf{میانگین} & \lr{۰٫۵۷۶} & \lr{۰٫۸۰۱} & \lr{۰٫۸۱۴} & \lr{۰٫۸۲۷} & \lr{۰٫۶۴۵} & \lr{۰٫۸۷۲} & \lr{۳۳٫۷} & \lr{۶۵٫۸} & \lr{۶۷٫۸} & \lr{۶۹٫۹} & \lr{۳۹٫۱} & \lr{۷۵٫۸} \\
\bottomrule
\end{tabular}
}
\end{table}


\fig{ctrl_A1_phase_vs_E_r.png}{0.95}{میانگین $r$ روی پنج عضله برای هر مدل. نقطه‌ی بزرگ میانگین ۱۲ نفر با خطای معیار است و هر خط خاکستری یک شرکت‌کننده. چهار مدل خاکستری فقط زمان رخدادها را می‌دانند؛ \lr{E} فقط \lr{EEG} و \lr{K} فقط سینماتیک.}

\paragraph{شکل \ref{fig:ctrl_A1_phase_vs_E_r.png} چه می‌گوید.} از چپ به راست، هر مدل اطلاعات بیشتری درباره‌ی زمان‌بندی دارد. فقط دانستن زمان روشن شدن \lr{LED} به $r$ \lr{۰٫۵۸} می‌رسد. با هشت رخداد این عدد به \lr{۰٫۸۰} و با جزئیات بیشتر به \lr{۰٫۸۱} و \lr{۰٫۸۳} می‌رسد. \lr{E} با \lr{۰٫۶۴} بین مدل \lr{cue} و مدل‌های زمان‌بندی کامل است. سه نکته از خط‌های افراد پیداست:
\begin{itemize}
\item مدل \lr{phase} در هر ۱۲ نفر از \lr{E} بهتر است. میانگین فاصله \lr{۰٫۱۷} واحد $r$ است و آزمون \lr{Wilcoxon} روی افراد $p=$ \lr{۰٫۰۰۰۵} می‌دهد.
\item در سه نفر (\lr{P1}، \lr{P3} و \lr{P9}) حتی مدل \lr{cue}، که فقط زمان روشن شدن چراغ را می‌داند، از \lr{E} بهتر است. با \lr{VAF} هم همین تصویر دیده می‌شود: \lr{cue} \lr{۳۳٫۷} درصد و \lr{E} \lr{۳۹٫۱} درصد (جدول \ref{tab:ctrl-phase}).
\item \lr{K} با \lr{۰٫۸۷} تنها مدلی است که از زمان‌بندی کامل بهتر است، و این در هر ۱۲ نفر صادق است.
\end{itemize}

\fig{ctrl_A3_phase_per_muscle.png}{0.95}{همان مقایسه برای هر عضله (میانگین ۱۲ نفر با خطای معیار).}

\paragraph{شکل \ref{fig:ctrl_A3_phase_per_muscle.png} چه می‌گوید.} این شکل یکی از نتیجه‌های مهم این فصل را دارد. فصل \ref{sec:literature} نشان داده بود که \lr{E} عضله‌ی شانه (\lr{AD}) را بهتر از همه و عضله‌های انگشت (\lr{FD} و \lr{FDI}) را بدتر از همه پیش‌بینی می‌کند، برعکس انتظار از \lr{CMC}. اینجا مدل زمان‌بندی، که هیچ اطلاعاتی از مغز ندارد، دقیقاً همین ترتیب را دارد: \lr{AD} \lr{۰٫۹۳}، \lr{ED} \lr{۰٫۹۰}، \lr{BR} \lr{۰٫۸۶}، \lr{FDI} \lr{۰٫۷۱} و \lr{FD} \lr{۰٫۶۷}. فاصله‌ی \lr{E} با مدل \lr{phase} هم در همه‌ی عضله‌ها تقریباً یکی است (\lr{۰٫۱۵} تا \lr{۰٫۲۰}). پس ترتیب عضله‌ها در \lr{E} از خود \lr{EEG} نمی‌آید. از این می‌آید که فعالیت \lr{AD} بیشتر از عضله‌های انگشت از روی مرحله‌ی تکلیف قابل حدس است.

\fig{ctrl_A4_phase_example.png}{0.95}{بیست ثانیه از اولین سری آزمون \lr{P1}. ناحیه‌ی خاکستری پوش واقعی \lr{EMG} است، خط‌چین مدل \lr{cue} و خط سیاه مدل \lr{phase}.}

\paragraph{شکل \ref{fig:ctrl_A4_phase_example.png} چه می‌گوید.} مدل \lr{phase} شکل هر بلند کردن را دنبال می‌کند: شروع فعالیت، سطح نگه‌داشتن، قله‌ی کوتاه \lr{FDI} هنگام گرفتن و پایان فعالیت. مدل \lr{cue} شروع را درست می‌گیرد، چون حرکت چند صد میلی‌ثانیه بعد از \lr{LED} شروع می‌شود، اما پایان را نمی‌داند و زود پایین می‌آید. آنچه هیچ‌کدام نمی‌گیرند نوسان‌های سریع درون هر بلند کردن و تفاوت دامنه بین بلند کردن‌هاست. این دقیقاً همان چیزی است که \lr{K} از سینماتیک می‌خواند.

\paragraph{تفاوت افراد.} همبستگی دقت \lr{E} و دقت مدل \lr{phase} بین ۱۲ نفر فقط \lr{۰٫۳۸} است. یعنی کسی که \lr{EMG}اش راحت‌تر از زمان‌بندی حدس زده می‌شود، لزوماً \lr{E} بهتری ندارد. نمونه‌ی روشن \lr{P9} است: مدل \lr{phase} برای او \lr{۰٫۸۵} است ولی \lr{E} فقط \lr{۰٫۴۸} (\lr{VAF} \lr{۴٫۵} درصد). پس اینکه \lr{E} چقدر از زمان‌بندی را بازیابی کند، به کیفیت \lr{EEG} هر فرد بستگی دارد.

\paragraph{نتیجه.} دقت \lr{E} در هیچ شرکت‌کننده‌ای بالاتر از چیزی نیست که دانستن زمان‌بندی تکلیف به دست می‌دهد، و ترتیب عضله‌ها در \lr{E} همان ترتیب مدل زمان‌بندی است. این ثابت نمی‌کند که \lr{E} فقط زمان‌بندی را می‌خواند. اما نشان می‌دهد که برای توضیح دقت \lr{E} به چیزی بیش از زمان‌بندی نیازی نیست. نکته‌ی مهم دیگر نتیجه‌ی مدل \lr{cue} است. پاسخ مغز به یک محرک بینایی، یا حرکت چشم به سمت \lr{LED}، کافی است تا بخش بزرگی از دقت \lr{E} به دست آید.

% ---------------------------------------------------------------------------------------------
\FloatBarrier
\subsection{آزمایش ۲: نقشه‌های مرجع \lr{ERD} و \lr{CMC}}
\label{sec:controls-ref}

\paragraph{\lr{ERD}.} وقتی قشر حرکتی فعال می‌شود، توان باندهای آلفا (۸ تا ۱۲ هرتز) و بتا (۱۳ تا ۳۰ هرتز) روی آن کم می‌شود \cite{pfurtscheller1999}. برای هر الکترود و هر بلند کردن:
\begin{equation}
\mathrm{ERD}\% = 100\cdot\frac{P_{\text{\lr{window}}}-P_{\text{\lr{baseline}}}}{P_{\text{\lr{baseline}}}},
\end{equation}
که $P$ توان باند است. \lr{baseline} بازه‌ی ۲ تا ۱ ثانیه پیش از شروع حرکت دست است، وقتی فرد آرام منتظر است. سه پنجره سنجیده شد: ثانیه‌ی پیش از اولین تماس، ثانیه‌ی پیش از \lr{lift-off}، و نگه‌داشتن (\lr{۰٫۵} تا \lr{۱٫۵} ثانیه بعد از \lr{lift-off}). عدد منفی یعنی کاهش توان، یعنی فعال شدن قشر زیر آن الکترود. نقشه‌ی هر فرد میانگین همه‌ی بلند کردن‌های همه‌ی سری‌های اوست.

\paragraph{\lr{CMC}.} هم‌دوسی بین \lr{EEG} هر الکترود و \lr{EMG} هر عضله در باند بتا (۱۵ تا ۳۰ هرتز)، در دوره‌ی نگه‌داشتن که نیروی عضله تقریباً ثابت است \cite{kilner2000}. \lr{EMG} از سیگنال خام ۴۰۰۰ هرتزی ساخته می‌شود (فیلتر ۳۰ تا ۳۰۰ هرتز، یکسوسازی، کاهش نرخ به ۵۰۰ هرتز) و هم‌دوسی روی قطعه‌های یک‌ثانیه‌ای با پنجره‌ی \lr{Hann} حساب می‌شود. سطح معناداری ۹۵ درصد برای یک نفر از تعداد قطعه‌ها به دست می‌آید و حدود \lr{۰٫۰۰۴} است.

\paragraph{مقایسه با \lr{E}.} برای هر نفر، سهم هر الکترود در \lr{rollout} و \lr{gradient-weighted rollout} خروجی \lr{E} (فصل \ref{sec:consistency}) با هر نقشه‌ی مرجع مقایسه می‌شود: همبستگی اسپیرمن روی ۳۲ الکترود. علامت \lr{ERD} برعکس می‌شود تا عدد مثبت یعنی «هر جا توان بیشتر کم شده، \lr{E} هم بیشتر از آن خوانده است». دو آزمون انجام می‌شود: \lr{sign-flip} روی ۱۲ نفر، و برای نقشه‌ی میانگین گروه، جابه‌جایی تصادفی برچسب الکترودها.

\fig{ctrl_B1_group_erd.png}{1.0}{نقشه‌ی گروهی \lr{ERD} (میانگین ۱۲ نفر). آبی یعنی کاهش توان نسبت به استراحت، قرمز یعنی افزایش.}

\paragraph{شکل \ref{fig:ctrl_B1_group_erd.png} چه می‌گوید.} در هر شش نقشه کانون کاهش توان روی ناحیه‌ی مرکزی دو نیمکره است، روی \lr{C3} و \lr{C4} و الکترودهای \lr{CP} کنار آن‌ها. این همان الگوی کتاب‌های درسی برای حرکت دست است. کاهش بتا تا نگه‌داشتن ادامه دارد و در آن مرحله به نظر می‌رسد بیشتر روی نیمکره‌ی چپ، یعنی نیمکره‌ی مقابل دست راست، متمرکز شود. پیشانی و پس‌سر تغییر کمتری دارند.

\begin{table}[htbp]
\centering
\footnotesize
\caption{هشت الکترود با بیشترین کاهش توان بتا در ثانیه‌ی پیش از \lr{lift-off}، میانگین گروه. عدد منفی یعنی توان نسبت به استراحت پیش از حرکت کم شده است. دو ستون آخر سهم همان الکترود در \lr{rollout} خروجی \lr{E} هنگام \lr{lift-off} (سهم یکنواخت \lr{۳٫۱۲۵} درصد) و رتبه‌ی آن بین ۳۲ الکترود.}
\label{tab:ctrl-erd-top}
\begin{tabular}{lccccc}
\toprule
\lr{electrode} & \lr{ERD} بتا (درصد) & \lr{ERD} آلفا (درصد) & افراد با کاهش بتا & سهم \lr{rollout} \lr{E} (درصد) & رتبه در \lr{rollout} \\
\midrule
\lr{C3} & \lr{−۴۱٫۳} & \lr{−۳۵٫۸} & \lr{۱۲/۱۲} & \lr{۵٫۵۳} & \lr{۱} \\
\lr{CP1} & \lr{−۳۹٫۸} & \lr{−۳۲٫۱} & \lr{۱۲/۱۲} & \lr{۴٫۹۸} & \lr{۲} \\
\lr{C4} & \lr{−۳۸٫۷} & \lr{−۴۱٫۶} & \lr{۱۲/۱۲} & \lr{۴٫۳۶} & \lr{۳} \\
\lr{CP2} & \lr{−۳۶٫۲} & \lr{−۳۳٫۷} & \lr{۱۲/۱۲} & \lr{۳٫۲۵} & \lr{۸} \\
\lr{P3} & \lr{−۳۵٫۹} & \lr{−۲۷٫۹} & \lr{۱۲/۱۲} & \lr{۳٫۲۰} & \lr{۱۰} \\
\lr{CP5} & \lr{−۳۳٫۹} & \lr{−۳۰٫۳} & \lr{۱۲/۱۲} & \lr{۲٫۹۳} & \lr{۱۸} \\
\lr{Pz} & \lr{−۳۲٫۸} & \lr{−۲۳٫۵} & \lr{۱۲/۱۲} & \lr{۲٫۵۰} & \lr{۳۰} \\
\lr{P4} & \lr{−۳۱٫۳} & \lr{−۲۷٫۰} & \lr{۱۲/۱۲} & \lr{۲٫۵۳} & \lr{۲۹} \\
\bottomrule
\end{tabular}
\end{table}


\paragraph{جدول \ref{tab:ctrl-erd-top} چه می‌گوید.} این الگو در تک‌تک افراد هست. هشت الکترود اول جدول در هر ۱۲ نفر کاهش بتا دارند. در ۱۱ نفر از ۱۲ نفر قوی‌ترین کاهش روی \lr{C3}، \lr{C4} یا \lr{CP1} است؛ فقط در \lr{P9} اول \lr{P3} است که آن هم کنار \lr{CP1} است. اندازه‌ی اثر بین افراد خیلی فرق می‌کند: کاهش بتا روی \lr{C3} از \lr{−۱۷} درصد در \lr{P6} تا \lr{−۷۳} درصد در \lr{P7}. مهم‌ترین نکته‌ی جدول دو ستون آخر است: سه الکترود اول \lr{ERD}، یعنی \lr{C3}، \lr{CP1} و \lr{C4}، دقیقاً همان سه الکترود اول \lr{rollout} خروجی \lr{E} هستند، با همان ترتیب. اما الکترودهای آهیانه‌ای مثل \lr{Pz} و \lr{P4} با وجود \lr{ERD} قوی، در \lr{E} سهم کمی دارند.

\fig{ctrl_B2_group_cmc.png}{1.0}{نقشه‌ی گروهی \lr{CMC} بتا در نگه‌داشتن، برای هر عضله. مقیاس رنگ هر نقشه جداست.}

\paragraph{شکل \ref{fig:ctrl_B2_group_cmc.png} چه می‌گوید.} \lr{CMC} در این داده ضعیف است. مقدار گروهی در همه‌ی الکترودها زیر سطح معناداری یک نفر است. در سطح فرد هم الکترود معنادار نادر است: \lr{P1} برای \lr{AD} ۹ الکترود و برای \lr{ED} ۴ الکترود دارد، \lr{P8} برای \lr{AD} و \lr{P7} برای \lr{FDI} هر کدام ۳ الکترود، و بیشتر افراد برای بیشتر عضله‌ها هیچ. با این حال، شکل نقشه‌ها دو الگوی متفاوت دارد:
\begin{itemize}
\item برای \lr{FDI} (انگشت اشاره) و \lr{ED} کانون روی \lr{C3} و \lr{CP1} است، یعنی قشر حرکتی نیمکره‌ی مقابل دست. این همان چیزی است که پژوهش‌های پیشین برای عضله‌های دست گزارش کرده‌اند.
\item برای \lr{AD} و \lr{BR} کانون روی پس‌سر است (\lr{PO9}، \lr{O1}، \lr{Oz} و \lr{PO10}) و برای \lr{FD} روی گیجگاه چپ و پس‌سر (\lr{T7} و \lr{PO9}). این جای غیرعادی‌ای برای \lr{CMC} است و احتمالاً از فعالیت عضله‌های گردن و پوست سر می‌آید که با حرکت بازو هم‌زمان‌اند.
\end{itemize}
پس \lr{CMC} این داده مرجع ضعیف و آلوده‌ای است و باید با احتیاط به آن تکیه کرد.

\begin{table}[htbp]
\centering
\footnotesize
\caption{هشت الکترود با بیشترین \lr{CMC} بتا (۱۵ تا ۳۰ هرتز) در نگه‌داشتن، میانگین پنج عضله و 12 نفر. سطح معناداری ۹۵ درصد برای یک نفر حدود \lr{۰٫۰۰۳۶} است.}
\label{tab:ctrl-cmc-top}
\begin{tabular}{lccc}
\toprule
\lr{electrode} & \lr{CMC} (میانگین عضله‌ها) & \lr{ERD} بتا (درصد) & سهم \lr{rollout} \lr{E} (درصد) \\
\midrule
\lr{PO9} & \lr{۰٫۰۰۱۹} & \lr{−۷٫۷} & \lr{۳٫۲۴} \\
\lr{O1} & \lr{۰٫۰۰۱۷} & \lr{−۲۰٫۹} & \lr{۳٫۴۳} \\
\lr{Oz} & \lr{۰٫۰۰۱۷} & \lr{−۱۸٫۵} & \lr{۲٫۹۵} \\
\lr{C3} & \lr{۰٫۰۰۱۷} & \lr{−۴۱٫۳} & \lr{۵٫۵۳} \\
\lr{PO10} & \lr{۰٫۰۰۱۶} & \lr{−۱۱٫۸} & \lr{۳٫۷۲} \\
\lr{O2} & \lr{۰٫۰۰۱۶} & \lr{−۱۷٫۱} & \lr{۳٫۰۷} \\
\lr{CP1} & \lr{۰٫۰۰۱۶} & \lr{−۳۹٫۸} & \lr{۴٫۹۸} \\
\lr{CP5} & \lr{۰٫۰۰۱۴} & \lr{−۳۳٫۹} & \lr{۲٫۹۳} \\
\bottomrule
\end{tabular}
\end{table}

\begin{table}[htbp]
\centering
\footnotesize
\caption{همبستگی اسپیرمن نقشه‌ی الکترود خروجی \lr{E} هنگام \lr{lift-off} با نقشه‌های مرجع، روی ۳۲ الکترود. پنج ستون اول \lr{rollout} و پنج ستون بعد \lr{gradient-weighted rollout}. $\bar\rho$: میانگین روی افراد؛ «مثبت»: شمار افراد با $\rho>0$؛ \lr{p} بعد از آن: \lr{sign-flip} روی افراد. $\rho$ گروه: نقشه‌ی میانگین گروه در برابر نقشه‌ی مرجع میانگین گروه؛ \lr{p} آن: جابه‌جایی برچسب الکترودها. علامت \lr{ERD} برعکس شده تا عدد مثبت یعنی «جایی که توان بیشتر کم شده، \lr{rollout} هم بیشتر است».}
\label{tab:ctrl-ref-rho}
\resizebox{\textwidth}{!}{%
\begin{tabular}{lcccccccccc}
\toprule
 & \multicolumn{5}{c}{\lr{rollout}} & \multicolumn{5}{c}{\lr{gradient-weighted rollout}} \\ \cmidrule(lr){2-6} \cmidrule(lr){7-11}
نقشه‌ی مرجع & $\bar\rho$ & مثبت & \lr{p} & $\rho$ گروه & \lr{p} & $\bar\rho$ & مثبت & \lr{p} & $\rho$ گروه & \lr{p} \\
\midrule
\lr{$-$ERD} آلفا، \lr{lift-off} & \lr{+۰٫۱۳} & \lr{۸/۱۲} & \lr{۰٫۰۶} & \lr{+۰٫۱۹} & \lr{۰٫۳۰۸} & \lr{+۰٫۱۱} & \lr{۹/۱۲} & \lr{۰٫۱۲} & \lr{+۰٫۱۳} & \lr{۰٫۴۸۶} \\
\lr{$-$ERD} بتا، \lr{lift-off} & \lr{+۰٫۲۲} & \lr{۹/۱۲} & \lr{۰٫۰۱} & \lr{+۰٫۲۸} & \lr{۰٫۱۱۴} & \lr{+۰٫۱۶} & \lr{۹/۱۲} & \lr{۰٫۰۴} & \lr{+۰٫۲۶} & \lr{۰٫۱۶۱} \\
\lr{$-$ERD} آلفا، نگه‌داشتن & \lr{+۰٫۱۴} & \lr{۱۰/۱۲} & \lr{۰٫۰۲} & \lr{+۰٫۳۷} & \lr{۰٫۰۴۰} & \lr{+۰٫۰۸} & \lr{۸/۱۲} & \lr{۰٫۱۳} & \lr{+۰٫۱۵} & \lr{۰٫۴۱۳} \\
\lr{$-$ERD} بتا، نگه‌داشتن & \lr{+۰٫۲۱} & \lr{۱۱/۱۲} & \lr{۰٫۰۱} & \lr{+۰٫۴۰} & \lr{۰٫۰۲۶} & \lr{+۰٫۱۱} & \lr{۸/۱۲} & \lr{۰٫۱۰} & \lr{+۰٫۲۲} & \lr{۰٫۲۳۲} \\
\lr{CMC} بتا، نگه‌داشتن & \lr{+۰٫۲۰} & \lr{۹/۱۲} & \lr{۰٫۰۰} & \lr{+۰٫۵۲} & \lr{۰٫۰۰۴} & \lr{+۰٫۲۰} & \lr{۹/۱۲} & \lr{۰٫۰۱} & \lr{+۰٫۴۵} & \lr{۰٫۰۱۱} \\
\bottomrule
\end{tabular}
}
\end{table}


\fig{ctrl_B3_rollout_vs_reference_lift_off.png}{0.95}{همبستگی اسپیرمن نقشه‌ی الکترود \lr{E} هنگام \lr{lift-off} با هر نقشه‌ی مرجع. نقطه‌های کوچک افرادند، نقطه‌ی بزرگ میانگین و خط بازه‌ی اطمینان ۹۵ درصد. آبی \lr{rollout} و سبز \lr{gradient-weighted rollout} است. عدد سمت راست $p$ آزمون \lr{sign-flip} است.}

\paragraph{شکل \ref{fig:ctrl_B3_rollout_vs_reference_lift_off.png} و جدول \ref{tab:ctrl-ref-rho} چه می‌گویند.} همه‌ی میانگین‌ها مثبت‌اند و بیشترشان در سطح گروه معنادارند:
\begin{itemize}
\item \textbf{با \lr{ERD} بتا:} برای \lr{rollout} میانگین $\rho$ \lr{۰٫۲۲} در ثانیه‌ی پیش از \lr{lift-off} ($p=$ \lr{۰٫۰۰۶}، مثبت در ۹ نفر) و \lr{۰٫۲۱} در نگه‌داشتن ($p=$ \lr{۰٫۰۰۹}، مثبت در ۱۱ نفر).
\item \textbf{با \lr{ERD} آلفا:} ضعیف‌تر، \lr{۰٫۱۳} و \lr{۰٫۱۴}، که فقط در نگه‌داشتن معنادار است.
\item \textbf{با \lr{CMC}:} میانگین \lr{۰٫۲۰} ($p=$ \lr{۰٫۰۰۴}). در نقشه‌ی میانگین گروه این عدد به \lr{۰٫۵۲} می‌رسد ($p=$ \lr{۰٫۰۰۴}).
\item \textbf{\lr{gradient-weighted rollout}} همین جهت را با اعداد کمی کوچک‌تر نشان می‌دهد؛ با \lr{ERD} بتا پیش از \lr{lift-off} و با \lr{CMC} هنوز معنادار است.
\end{itemize}
اندازه‌ی این همبستگی‌ها متوسط است: $\rho=0.2$ یعنی ترتیب الکترودها در دو نقشه فقط کمی شبیه است. افراد هم با هم فرق دارند و بسته به نقشه‌ی مرجع، در یک تا سه نفر همبستگی منفی است. شکل \lr{first-touch} (\lr{ctrl\_B4} در پوشه‌ی شکل‌ها) همین الگو را دارد: با \lr{ERD} بتا \lr{۰٫۲۱} ($p=$ \lr{۰٫۰۰۶}) و با \lr{CMC} \lr{۰٫۲۰} ($p=$ \lr{۰٫۰۰۶}).

\fig{ctrl_B5_electrode_scatter.png}{0.95}{هر نقطه یک الکترود است. محور عمودی سهم الکترود در \lr{rollout} گروهی \lr{E} و محور افقی نقشه‌ی مرجع گروهی است: چپ، کاهش بتا پیش از \lr{lift-off}؛ راست، \lr{CMC} بتا (میانگین عضله‌ها).}

\paragraph{شکل \ref{fig:ctrl_B5_electrode_scatter.png} چه می‌گوید.} این شکل نشان می‌دهد چرا همبستگی‌ها متوسط‌اند. الکترودها در نمودار چپ سه گروه می‌سازند:
\begin{itemize}
\item \textbf{\lr{C3}، \lr{CP1} و \lr{C4}:} بالا و سمت راست. هم \lr{E} آن‌ها را بیشتر از همه می‌خواند و هم بیشترین \lr{ERD} را دارند. این بخش با قشر حسی-حرکتی می‌خواند و بزرگ‌ترین سهم \lr{E} است.
\item \textbf{\lr{PO10}، \lr{PO9}، \lr{O1}، \lr{F7} و \lr{T7}:} بالا و سمت چپ. \lr{E} آن‌ها را بیشتر از سهم یکنواخت می‌خواند، اما \lr{ERD} کمی دارند. \lr{F7} کنار چشم است و حرکت چشم را می‌گیرد؛ \lr{PO9}، \lr{PO10} و \lr{O1} روی عضله‌های پشت گردن‌اند و \lr{T7} روی عضله‌ی گیجگاهی.
\item \textbf{\lr{Pz}، \lr{P4} و \lr{CP5}:} پایین و سمت راست. \lr{ERD} قوی دارند ولی \lr{E} کمتر از سهم یکنواخت از آن‌ها می‌خواند.
\end{itemize}
در نمودار راست همبستگی با \lr{CMC} بیشتر است ($\rho=0.52$)، چون \lr{CMC} هم در هر دو گروه اول بالاست: \lr{C3} و \lr{CP1} از یک طرف، و \lr{PO9}، \lr{O1} و \lr{PO10} از طرف دیگر. پس بخشی از این توافق از قشر حرکتی می‌آید و بخشی از سیگنال عضلانی مشترکی که هم \lr{E} و هم \lr{CMC} می‌بینند.

\paragraph{نتیجه.} نقشه‌ی \lr{E} در سطح گروه به‌طور معنادار با \lr{ERD} و \lr{CMC} هم‌جهت است و کانون اصلی‌اش همان کانون \lr{ERD} است: \lr{C3}، \lr{CP1} و \lr{C4}. این قوی‌ترین شاهد این گزارش است که \lr{E} دست‌کم بخشی از اطلاعاتش را از قشر حسی-حرکتی می‌گیرد. اما دو چیز تصویر را کامل نمی‌کند. اول، \lr{E} از الکترودهای کنار چشم، گیجگاه و پشت گردن هم می‌خواند، که جای معمول \lr{artifact}اند. دوم، اندازه‌ی همبستگی متوسط است. آزمایش بعدی باید \lr{E} را یک بار فقط با الکترودهای مرکزی و یک بار فقط با الکترودهای کناری آموزش دهد تا سهم هر گروه مستقیم اندازه گرفته شود.

% ---------------------------------------------------------------------------------------------
\FloatBarrier
\subsection{آزمایش ۳: پیش‌بینی آینده با مدل علّی}
\label{sec:controls-horizon}

\paragraph{ایده.} سینماتیک نتیجه‌ی فعالیت عضله است و کمی بعد از آن می‌آید. فرمان قشر حرکتی پیش از هر دو است: \lr{ERD} حدود یک تا دو ثانیه پیش از حرکت شروع می‌شود. اگر \lr{EEG} اطلاعاتی داشته باشد که سینماتیک ندارد، این اطلاعات باید بیشتر درباره‌ی \emph{آینده} باشد. پس مدل را مجبور می‌کنیم \lr{EMG} را $h$ میلی‌ثانیه جلوتر پیش‌بینی کند: ورودی‌ها تا لحظه‌ی $t$، هدف در لحظه‌ی $t+h$. پیش‌بینی ما این بود که هر چه $h$ بزرگ‌تر شود، سهم \lr{EEG} بیشتر شود، به‌ویژه درست پیش از شروع حرکت که سینماتیک هنوز چیزی نمی‌بیند.

\paragraph{مدل علّی.} این آزمایش فقط وقتی معنا دارد که مدل واقعاً به آینده دسترسی نداشته باشد. \lr{v2} برون‌خط است: هر \lr{token} کل پنجره‌ی ۸ ثانیه‌ای را می‌بیند. برای این آزمایش یک حالت علّی به \lr{v2} اضافه شد که با یک کلید فعال می‌شود و رفتار پیش‌فرض را تغییر نمی‌دهد:
\begin{itemize}
\item \textbf{مدل:} \lr{attention} در هر دو \lr{encoder} و در \lr{fusion} یک \lr{causal mask} می‌گیرد، یعنی هر \lr{token} فقط \lr{token}های گذشته و خودش را می‌بیند. \lr{convolution} ورودی سینماتیک فقط از سمت چپ \lr{padding} می‌گیرد.
\item \textbf{پیش‌پردازش ورودی‌ها:} فیلترهای صفرفاز (\lr{filtfilt}) که از آینده هم استفاده می‌کنند، با فیلترهای یک‌طرفه جایگزین شدند: \lr{Butterworth} مرتبه‌ی ۴ با \lr{sosfilt} برای کاهش نرخ، \lr{band power} با فیلتر میان‌گذر علّی و دو \lr{low-pass} مرتبه‌ی اول، و سرعت با مشتق \lr{Savitzky-Golay} یک‌طرفه.
\item \textbf{پیش‌بینی پیوسته:} از هر پنجره، انتهای آن نگه داشته می‌شود (نه وسطش) تا هر \lr{token} دست‌کم ۴ ثانیه گذشته داشته باشد.
\end{itemize}
آزمون‌های واحد بررسی می‌کنند که هیچ خروجی به ورودی بعد از خودش وابسته نیست. یک نشت کوچک باقی است: \lr{gate} نیرو در ویژگی‌های تاریخچه با میانگین متحرک وسط‌چین ساخته می‌شود و تا حدود ۵۰ میلی‌ثانیه آینده را می‌بیند.

\paragraph{سنجه‌ها.} برای هر افق ($h$ برابر صفر، ۲۰۰، ۵۰۰ و ۱۰۰۰ میلی‌ثانیه) یک مدل جدا برای هر نفر آموزش دید، یعنی ۴۸ مدل. سه سنجه سهم \lr{EEG} را از سه زاویه نشان می‌دهند:
\begin{itemize}
\item $F-K$: بهره‌ی \lr{VAF} خروجی ترکیبی نسبت به سینماتیک تنها؛
\item $r$ \lr{residual test} (فصل \ref{sec:residual}): آیا \lr{EEG} خطای \lr{K} را پیش‌بینی می‌کند؟
\item همین آزمون جدا برای هر مرحله‌ی تکلیف، با مهارت
\begin{equation}
\text{\lr{skill}}=1-\frac{\mathrm{MSE}(K+\hat\varepsilon)}{\mathrm{MSE}(K)},
\end{equation}
که $\hat\varepsilon$ خطای پیش‌بینی‌شده‌ی \lr{K} از \lr{EEG} است. مهارت مثبت یعنی اضافه کردن \lr{EEG} خطا را کم می‌کند. مرحله‌ی «پیش از حرکت» ثانیه‌ی پیش از شروع حرکت دست است. معناداری برای هر فرد در برابر همان پیش‌بینی با جابه‌جایی زمانی سنجیده می‌شود.
\end{itemize}

\begin{table}[htbp]
\centering
\footnotesize
\caption{مدل علّی با افق‌های پیش‌بینی مختلف، میانگین 12 نفر ($\pm$ خطای معیار). ستون‌های \lr{p}: آزمون \lr{Wilcoxon} یک‌طرفه روی افراد، برای اینکه \lr{F}$-$\lr{K} یا مهارت بزرگ‌تر از صفر باشد. $r$ باقی‌مانده: \lr{residual test} با \lr{tokens} \lr{EEG} (میانگین عضله‌ها)؛ «معنادار»: شمار افراد با $p\le0.05$. مهارت پیش از حرکت: $1-\mathrm{MSE}(K+\mathrm{EEG})/\mathrm{MSE}(K)$ در ثانیه‌ی پیش از شروع حرکت دست؛ «معنادار»: مهارت مثبت و $p\le0.05$ در برابر جابه‌جایی زمانی.}
\label{tab:ctrl-horizon}
\resizebox{\textwidth}{!}{%
\begin{tabular}{lccccccccccc}
\toprule
افق (\lr{ms}) & \lr{VAF} \lr{K} & $r$ \lr{K} & $r$ \lr{E} & \lr{F}$-$\lr{K} (\lr{VAF}) & \lr{p} & $r$ باقی‌مانده & معنادار & مهارت \lr{ridge} پیش از حرکت & معنادار & \lr{p} & مهارت \lr{F} پیش از حرکت \\
\midrule
\lr{۰} & \lr{۷۵٫۵}\,$\pm$\,\lr{۲٫۶} & \lr{۰٫۸۷۲} & \lr{۰٫۵۶۲} & \lr{+۰٫۰۸}\,$\pm$\,\lr{۰٫۰۴} & \lr{۰٫۰۳۹} & \lr{۰٫۰۸۳} & \lr{۱۱/۱۲} & \lr{−۰٫۰۰۳} & \lr{۶/۱۲} & \lr{۰٫۱۹} & \lr{+۰٫۰۱۳} \\
\lr{۲۰۰} & \lr{۷۰٫۹}\,$\pm$\,\lr{۲٫۹} & \lr{۰٫۸۴۶} & \lr{۰٫۵۴۹} & \lr{+۰٫۲۴}\,$\pm$\,\lr{۰٫۱۰} & \lr{۰٫۰۴۶} & \lr{۰٫۰۹۲} & \lr{۱۱/۱۲} & \lr{−۰٫۰۰۲} & \lr{۶/۱۲} & \lr{۰٫۲۳} & \lr{+۰٫۰۰۳} \\
\lr{۵۰۰} & \lr{۶۶٫۲}\,$\pm$\,\lr{۳٫۰} & \lr{۰٫۸۱۸} & \lr{۰٫۵۳۰} & \lr{+۰٫۵۵}\,$\pm$\,\lr{۰٫۱۶} & \lr{۰٫۰۰۰} & \lr{۰٫۰۸۹} & \lr{۱۲/۱۲} & \lr{−۰٫۰۴۰} & \lr{۲/۱۲} & \lr{۱٫۰۰} & \lr{−۰٫۰۱۷} \\
\lr{۱۰۰۰} & \lr{۵۵٫۵}\,$\pm$\,\lr{۳٫۷} & \lr{۰٫۷۵۹} & \lr{۰٫۴۳۹} & \lr{+۰٫۴۰}\,$\pm$\,\lr{۰٫۳۰} & \lr{۰٫۱۰۲} & \lr{۰٫۰۹۵} & \lr{۹/۱۲} & \lr{−۰٫۰۳۲} & \lr{۳/۱۲} & \lr{۰٫۹۷} & \lr{+۰٫۰۱۲} \\
\bottomrule
\end{tabular}
}
\end{table}


\fig{ctrl_C2_accuracy_vs_horizon.png}{0.95}{دقت مدل علّی در برابر افق پیش‌بینی. چپ: \lr{VAF} خروجی \lr{K}؛ راست: $r$ خروجی \lr{E}. خط‌های خاکستری افرادند.}

\paragraph{شکل \ref{fig:ctrl_C2_accuracy_vs_horizon.png} چه می‌گوید.} اول، مدل علّی در $h=0$ برای \lr{K} تقریباً چیزی از دست نمی‌دهد: \lr{K} علّی به \lr{VAF} \lr{۷۵٫۵} درصد می‌رسد و \lr{K} برون‌خط \lr{۷۵٫۸} درصد بود (جدول \ref{tab:v2-subj}). \lr{E} علّی ضعیف‌تر است ($r$ \lr{۰٫۵۶} در برابر \lr{۰٫۶۴}). پس \lr{E} برون‌خط بخشی از دقتش را از دیدن \lr{EEG} بعد از هر لحظه می‌گرفت. دوم، با بزرگ شدن افق \lr{K} سریع افت می‌کند: از \lr{۷۵٫۵} به \lr{۵۵٫۵} درصد در یک ثانیه. \lr{E} آهسته‌تر افت می‌کند: $r$ آن تا ۵۰۰ میلی‌ثانیه تقریباً ثابت است و در یک ثانیه به \lr{۰٫۴۴} می‌رسد. به زبان نسبی، \lr{K} در یک ثانیه حدود ۲۶ درصد از \lr{VAF} خود را از دست می‌دهد و \lr{E} حدود ۲۲ درصد از $r$ خود را. این همان جهتی است که پیش‌بینی کرده بودیم، ولی تفاوت کوچک است.

\fig{ctrl_C1_eeg_gain_vs_horizon.png}{0.95}{سهم \lr{EEG} در برابر افق. چپ: $F-K$ (واحد \lr{VAF})؛ وسط: $r$ \lr{residual test} روی کل سری‌ها؛ راست: مهارت \lr{ridge} در ثانیه‌ی پیش از حرکت. نقطه‌ی آبی میانگین با خطای معیار و خط‌های خاکستری افرادند.}

\paragraph{شکل \ref{fig:ctrl_C1_eeg_gain_vs_horizon.png} چه می‌گوید.} سه سنجه سه جواب متفاوت می‌دهند:
\begin{itemize}
\item \textbf{$F-K$ با افق رشد می‌کند.} از \lr{+۰٫۰۸} در $h=0$ به \lr{+۰٫۲۴} و \lr{+۰٫۵۵} در ۲۰۰ و ۵۰۰ میلی‌ثانیه می‌رسد. در ۵۰۰ میلی‌ثانیه در ۱۱ نفر از ۱۲ نفر مثبت است و در سطح گروه معنادار ($p=$ \lr{۰٫۰۰۰۵}). در یک ثانیه میانگین \lr{+۰٫۴۰} است ولی دو نفر (\lr{P4} و \lr{P5}) حدود \lr{−۱٫۵} واحد افت دارند و آزمون گروهی معنادار نیست. بزرگ‌ترین بهره‌ها حدود ۱ تا ۲ واحد \lr{VAF} است (\lr{P3}، \lr{P8}، \lr{P10} و \lr{P11})، جدول \ref{tab:ctrl-horizon-subj}.
\item \textbf{$r$ باقی‌مانده تقریباً ثابت است.} \lr{۰٫۰۸۳}، \lr{۰٫۰۹۲}، \lr{۰٫۰۸۹} و \lr{۰٫۰۹۵}، و در ۹ تا ۱۲ نفر معنادار. اما همین آزمون با ویژگی‌های خام \lr{band power} با افق \emph{کم} می‌شود: از \lr{۰٫۰۴۰} (معنادار در ۸ نفر) به \lr{۰٫۰۲۲} (معنادار در ۲ نفر).
\item \textbf{مهارت پیش از حرکت مثبت نمی‌شود.} \lr{−۰٫۰۰۳}، \lr{−۰٫۰۰۲}، \lr{−۰٫۰۴۰} و \lr{−۰٫۰۳۲}. در هیچ افقی آزمون گروهی آن را بزرگ‌تر از صفر نشان نمی‌دهد. یعنی اضافه کردن پیش‌بینی \lr{EEG} به \lr{K} خطای ثانیه‌ی پیش از حرکت را کم نمی‌کند.
\end{itemize}

\begin{table}[htbp]
\centering
\footnotesize
\caption{\lr{residual test} جدا برای هر مرحله‌ی تکلیف و هر افق، میانگین 12 نفر. $r$: همبستگی باقی‌مانده‌ی پیش‌بینی‌شده‌ی \lr{ridge} با باقی‌مانده‌ی واقعی \lr{K} در همان مرحله. مهارت: $1-\mathrm{MSE}(K+\mathrm{EEG})/\mathrm{MSE}(K)$؛ «معنادار»: شمار افراد با مهارت مثبت و $p\le0.05$ در برابر جابه‌جایی زمانی. دو ستون آخر: شمار افرادی که $r$ آن‌ها از افق صفر تا بلندترین افق بالا رفته و \lr{p} آزمون \lr{Wilcoxon} یک‌طرفه برای این افزایش.}
\label{tab:ctrl-segments}
\resizebox{\textwidth}{!}{%
\begin{tabular}{lcccccccccccccc}
\toprule
 & \multicolumn{3}{c}{افق \lr{۰} میلی‌ثانیه} & \multicolumn{3}{c}{افق \lr{۲۰۰} میلی‌ثانیه} & \multicolumn{3}{c}{افق \lr{۵۰۰} میلی‌ثانیه} & \multicolumn{3}{c}{افق \lr{۱۰۰۰} میلی‌ثانیه} & \multicolumn{2}{c}{رشد $r$} \\ \cmidrule(lr){2-4} \cmidrule(lr){5-7} \cmidrule(lr){8-10} \cmidrule(lr){11-13} \cmidrule(lr){14-15}
مرحله & $r$ & مهارت & معنادار & $r$ & مهارت & معنادار & $r$ & مهارت & معنادار & $r$ & مهارت & معنادار & افزایش & \lr{p} \\
\midrule
پیش از حرکت & \lr{۰٫۰۶۲} & \lr{−۰٫۰۰۳} & \lr{۶/۱۲} & \lr{۰٫۱۰۰} & \lr{−۰٫۰۰۲} & \lr{۶/۱۲} & \lr{۰٫۱۰۹} & \lr{−۰٫۰۴۰} & \lr{۲/۱۲} & \lr{۰٫۱۹۰} & \lr{−۰٫۰۳۲} & \lr{۳/۱۲} & \lr{۱۰/۱۲} & \lr{۰٫۰۰۸} \\
رساندن دست & \lr{۰٫۰۷۱} & \lr{−۰٫۰۴۶} & \lr{۳/۱۲} & \lr{۰٫۱۸۲} & \lr{+۰٫۰۲۲} & \lr{۹/۱۲} & \lr{۰٫۱۷۲} & \lr{+۰٫۰۵۰} & \lr{۸/۱۲} & \lr{۰٫۱۶۲} & \lr{+۰٫۰۳۶} & \lr{۴/۱۲} & \lr{۹/۱۲} & \lr{۰٫۰۳۲} \\
بارگذاری & \lr{۰٫۰۱۷} & \lr{−۰٫۰۳۸} & \lr{۳/۱۲} & \lr{−۰٫۰۰۵} & \lr{−۰٫۰۳۳} & \lr{۳/۱۲} & \lr{۰٫۰۰۵} & \lr{−۰٫۰۱۹} & \lr{۲/۱۲} & \lr{۰٫۰۲۳} & \lr{−۰٫۰۱۹} & \lr{۰/۱۲} & \lr{۷/۱۲} & \lr{۰٫۳۹۶} \\
نگه‌داشتن & \lr{۰٫۰۲۴} & \lr{−۰٫۰۳۶} & \lr{۳/۱۲} & \lr{۰٫۰۱۱} & \lr{−۰٫۰۲۶} & \lr{۴/۱۲} & \lr{۰٫۰۱۴} & \lr{−۰٫۰۳۴} & \lr{۲/۱۲} & \lr{۰٫۰۲۷} & \lr{−۰٫۰۳۵} & \lr{۳/۱۲} & \lr{۷/۱۲} & \lr{۰٫۳۳۹} \\
رها کردن & \lr{۰٫۰۲۳} & \lr{−۰٫۰۲۶} & \lr{۲/۱۲} & \lr{۰٫۰۲۶} & \lr{−۰٫۰۱۶} & \lr{۳/۱۲} & \lr{۰٫۰۲۵} & \lr{−۰٫۰۲۱} & \lr{۱/۱۲} & \lr{۰٫۰۳۵} & \lr{−۰٫۰۲۳} & \lr{۱/۱۲} & \lr{۷/۱۲} & \lr{۰٫۲۱۲} \\
بقیه & \lr{۰٫۰۴۵} & \lr{−۰٫۰۷۴} & \lr{۴/۱۲} & \lr{۰٫۰۳۶} & \lr{−۰٫۰۵۹} & \lr{۴/۱۲} & \lr{۰٫۰۵۲} & \lr{−۰٫۰۷۲} & \lr{۱/۱۲} & \lr{۰٫۱۰۷} & \lr{−۰٫۰۱۲} & \lr{۵/۱۲} & \lr{۹/۱۲} & \lr{۰٫۰۳۲} \\
همه & \lr{۰٫۰۸۳} & \lr{−۰٫۰۴۰} & \lr{۴/۱۲} & \lr{۰٫۰۹۲} & \lr{−۰٫۰۱۹} & \lr{۵/۱۲} & \lr{۰٫۰۸۹} & \lr{−۰٫۰۱۸} & \lr{۴/۱۲} & \lr{۰٫۰۹۵} & \lr{−۰٫۰۰۹} & \lr{۵/۱۲} & \lr{۶/۱۲} & \lr{۰٫۳۹۶} \\
\bottomrule
\end{tabular}
}
\end{table}


\fig{ctrl_C3_gain_per_phase_ridge.png}{0.95}{مهارت \lr{ridge} برای هر مرحله‌ی تکلیف و هر افق، میانگین ۱۲ نفر. قرمز یعنی \lr{EEG} خطای \lr{K} را کم می‌کند و آبی یعنی زیاد. زیر هر عدد، شمار افرادی آمده که مهارتشان مثبت و معنادار است.}

\paragraph{شکل \ref{fig:ctrl_C3_gain_per_phase_ridge.png} و جدول \ref{tab:ctrl-segments} چه می‌گویند.} جدول دو چیز را جدا نشان می‌دهد: $r$، یعنی آیا پیش‌بینی \lr{EEG} جهت درست خطای \lr{K} را دارد، و مهارت، یعنی آیا آن پیش‌بینی خطا را کم هم می‌کند.
\begin{itemize}
\item \textbf{$r$ پیش از حرکت با افق رشد می‌کند.} \lr{۰٫۰۶۲}، \lr{۰٫۱۰۰}، \lr{۰٫۱۰۹} و \lr{۰٫۱۹۰}. در ۱۰ نفر از ۱۲ نفر $r$ در یک ثانیه از $r$ در افق صفر بیشتر است و این رشد در سطح گروه معنادار است ($p=$ \lr{۰٫۰۰۸}). در رساندن دست هم $r$ از \lr{۰٫۰۷۱} به حدود \lr{۰٫۱۷} می‌رسد ($p=$ \lr{۰٫۰۳}). در بارگذاری، نگه‌داشتن و رها کردن $r$ کوچک می‌ماند (\lr{۰٫۰۰} تا \lr{۰٫۰۴}) و رشد معناداری ندارد. این دقیقاً الگویی است که پیش‌بینی کرده بودیم: اطلاعات \lr{EEG} درباره‌ی خطای \lr{K} فقط در مرحله‌هایی رشد می‌کند که حرکت هنوز شروع نشده یا تازه شروع شده است.
\item \textbf{مهارت فقط در رساندن دست مثبت می‌شود.} \lr{+۰٫۰۲۲}، \lr{+۰٫۰۵۰} و \lr{+۰٫۰۳۶} در ۲۰۰، ۵۰۰ و ۱۰۰۰ میلی‌ثانیه، معنادار در ۹، ۸ و ۴ نفر از ۱۲ و در سطح گروه معنادار در هر سه افق. یعنی تا ۵ درصد کاهش خطای \lr{K} در این مرحله.
\item \textbf{ولی بخشی از این بهره به زمان بستگی ندارد.} میانگین \lr{null} برای رساندن دست در ۵۰۰ میلی‌ثانیه خودش \lr{+۰٫۰۲۶} است. یعنی همان پیش‌بینی، وقتی در زمان جابه‌جا شود، هنوز نیمی از بهره را می‌دهد. این بخش احتمالاً یک تصحیح ثابت است: \lr{K} در افق بلند فعالیت رساندن دست را کمتر از واقع برآورد می‌کند و هر پیش‌بینی مثبتی آن را کمی جبران می‌کند. بخش وابسته به زمان حدود \lr{۰٫۰۲} تا \lr{۰٫۰۳} است.
\item \textbf{پیش از حرکت، $r$ بالا ولی مهارت منفی است.} \lr{ridge} روی همه‌ی مرحله‌ها با هم برازش شده و اندازه‌ی اصلاحش برای این مرحله مناسب نیست: جهت درست است، ولی اندازه نه. میانگین \lr{null} هم منفی است (\lr{−۰٫۰۳} تا \lr{−۰٫۰۶})، چون هر اصلاحی که زمانش درست نباشد خطا را زیاد می‌کند. پس «معنادار» در این جدول یعنی «بهتر از اصلاحی هم‌اندازه با زمان تصادفی»، نه لزوماً بهتر از \lr{K}.
\end{itemize}

\fig{ctrl_C4_gain_per_phase_F.png}{0.95}{همان نمودار برای خروجی \lr{F} خود مدل (به جای \lr{ridge} بیرونی).}

\paragraph{شکل \ref{fig:ctrl_C4_gain_per_phase_F.png} چه می‌گوید.} اصلاح خود مدل در ۲۵ خانه از ۲۸ خانه کوچک و مثبت است (\lr{+۰٫۰۰۱} تا \lr{+۰٫۰۳۵}). برخلاف \lr{ridge}، \lr{fusion} یاد گرفته در مرحله‌هایی که کمکی از \lr{EEG} نمی‌آید تقریباً هیچ اصلاحی نکند. بیشترین بهره در رساندن دست و بارگذاری و در افق‌های ۵۰۰ و ۱۰۰۰ میلی‌ثانیه است (حدود \lr{+۰٫۰۳۵}). تنها خانه‌ی منفی مهم مرحله‌ی «بقیه» (بین بلند کردن‌ها) در یک ثانیه است (\lr{−۰٫۰۳۸})، که همان افت \lr{P4} و \lr{P5} در جدول \ref{tab:ctrl-horizon-subj} است.

\begin{table}[htbp]
\centering
\footnotesize
\caption{جزئیات جدول \ref{tab:ctrl-horizon} برای هر شرکت‌کننده. ستاره: $r$ باقی‌مانده با $p\le0.05$ در برابر جابه‌جایی زمانی.}
\label{tab:ctrl-horizon-subj}
\resizebox{\textwidth}{!}{%
\begin{tabular}{lcccccccccccc}
\toprule
 & \multicolumn{3}{c}{افق \lr{۰} میلی‌ثانیه} & \multicolumn{3}{c}{افق \lr{۲۰۰} میلی‌ثانیه} & \multicolumn{3}{c}{افق \lr{۵۰۰} میلی‌ثانیه} & \multicolumn{3}{c}{افق \lr{۱۰۰۰} میلی‌ثانیه} \\ \cmidrule(lr){2-4} \cmidrule(lr){5-7} \cmidrule(lr){8-10} \cmidrule(lr){11-13}
شرکت‌کننده & \lr{VAF} \lr{K} & \lr{F}$-$\lr{K} & $r$ باقی‌مانده & \lr{VAF} \lr{K} & \lr{F}$-$\lr{K} & $r$ باقی‌مانده & \lr{VAF} \lr{K} & \lr{F}$-$\lr{K} & $r$ باقی‌مانده & \lr{VAF} \lr{K} & \lr{F}$-$\lr{K} & $r$ باقی‌مانده \\
\midrule
\lr{P1} & \lr{۸۴٫۵} & \lr{−۰٫۰۰} & \lr{۰٫۰۷۱}$^{*}$ & \lr{۷۸٫۰} & \lr{+۰٫۰۱} & \lr{۰٫۰۷۸}$^{*}$ & \lr{۷۳٫۳} & \lr{+۰٫۰۶} & \lr{۰٫۰۵۶}$^{*}$ & \lr{۶۰٫۷} & \lr{+۰٫۱۰} & \lr{۰٫۰۷۰}$^{*}$ \\
\lr{P2} & \lr{۷۸٫۶} & \lr{+۰٫۰۲} & \lr{۰٫۰۷۳}$^{*}$ & \lr{۷۶٫۰} & \lr{−۰٫۰۳} & \lr{۰٫۰۸۰}$^{*}$ & \lr{۷۱٫۸} & \lr{+۰٫۰۹} & \lr{۰٫۱۴۹}$^{*}$ & \lr{۶۵٫۲} & \lr{+۰٫۱۰} & \lr{۰٫۰۱۱} \\
\lr{P3} & \lr{۷۰٫۳} & \lr{−۰٫۰۱} & \lr{۰٫۰۶۰}$^{*}$ & \lr{۶۲٫۵} & \lr{+۰٫۸۱} & \lr{۰٫۰۵۶}$^{*}$ & \lr{۵۵٫۷} & \lr{+۱٫۵۰} & \lr{۰٫۰۸۹}$^{*}$ & \lr{۳۵٫۷} & \lr{+۱٫۰۰} & \lr{۰٫۱۶۵}$^{*}$ \\
\lr{P4} & \lr{۸۱٫۸} & \lr{+۰٫۱۳} & \lr{۰٫۰۸۶}$^{*}$ & \lr{۷۷٫۲} & \lr{+۰٫۰۶} & \lr{۰٫۰۸۳}$^{*}$ & \lr{۶۸٫۶} & \lr{+۰٫۳۸} & \lr{۰٫۰۵۰}$^{*}$ & \lr{۵۱٫۵} & \lr{−۱٫۴۴} & \lr{۰٫۱۵۳}$^{*}$ \\
\lr{P5} & \lr{۶۳٫۰} & \lr{+۰٫۲۷} & \lr{۰٫۱۰۹}$^{*}$ & \lr{۵۷٫۶} & \lr{+۰٫۷۶} & \lr{۰٫۰۹۴}$^{*}$ & \lr{۵۰٫۷} & \lr{+۱٫۲۳} & \lr{۰٫۰۹۶}$^{*}$ & \lr{۳۹٫۹} & \lr{−۱٫۴۹} & \lr{۰٫۰۷۱}$^{*}$ \\
\lr{P6} & \lr{۷۶٫۹} & \lr{+۰٫۱۸} & \lr{۰٫۰۵۰}$^{*}$ & \lr{۷۴٫۰} & \lr{+۰٫۱۱} & \lr{۰٫۰۷۰}$^{*}$ & \lr{۷۰٫۲} & \lr{+۰٫۳۶} & \lr{۰٫۰۱۹}$^{*}$ & \lr{۵۷٫۱} & \lr{+۰٫۷۳} & \lr{۰٫۰۱۷} \\
\lr{P7} & \lr{۸۳٫۳} & \lr{−۰٫۰۵} & \lr{−۰٫۰۴۲} & \lr{۸۰٫۹} & \lr{−۰٫۰۶} & \lr{−۰٫۰۰۴} & \lr{۷۶٫۱} & \lr{+۰٫۱۸} & \lr{۰٫۰۳۶}$^{*}$ & \lr{۶۶٫۶} & \lr{+۰٫۷۱} & \lr{۰٫۱۴۷}$^{*}$ \\
\lr{P8} & \lr{۵۵٫۲} & \lr{+۰٫۰۳} & \lr{۰٫۲۶۱}$^{*}$ & \lr{۴۷٫۷} & \lr{−۰٫۱۷} & \lr{۰٫۲۰۸}$^{*}$ & \lr{۴۴٫۰} & \lr{+۱٫۵۲} & \lr{۰٫۲۳۳}$^{*}$ & \lr{۳۲٫۳} & \lr{+۱٫۹۶} & \lr{۰٫۱۰۱}$^{*}$ \\
\lr{P9} & \lr{۸۳٫۲} & \lr{+۰٫۰۲} & \lr{۰٫۰۷۰}$^{*}$ & \lr{۸۰٫۳} & \lr{−۰٫۰۲} & \lr{۰٫۰۶۵}$^{*}$ & \lr{۷۶٫۴} & \lr{−۰٫۰۲} & \lr{۰٫۰۸۸}$^{*}$ & \lr{۶۹٫۳} & \lr{−۰٫۰۱} & \lr{−۰٫۰۰۳} \\
\lr{P10} & \lr{۷۹٫۶} & \lr{+۰٫۱۲} & \lr{۰٫۰۲۵}$^{*}$ & \lr{۷۴٫۱} & \lr{+۰٫۷۷} & \lr{۰٫۱۰۳}$^{*}$ & \lr{۷۱٫۰} & \lr{+۰٫۳۵} & \lr{۰٫۰۷۱}$^{*}$ & \lr{۶۴٫۳} & \lr{+۱٫۱۵} & \lr{۰٫۱۶۶}$^{*}$ \\
\lr{P11} & \lr{۷۷٫۷} & \lr{+۰٫۴۰} & \lr{۰٫۱۰۹}$^{*}$ & \lr{۷۵٫۳} & \lr{+۰٫۲۳} & \lr{۰٫۱۲۸}$^{*}$ & \lr{۷۱٫۷} & \lr{+۰٫۷۵} & \lr{۰٫۰۹۱}$^{*}$ & \lr{۶۴٫۰} & \lr{+۱٫۲۲} & \lr{۰٫۱۱۵}$^{*}$ \\
\lr{P12} & \lr{۷۱٫۶} & \lr{−۰٫۱۰} & \lr{۰٫۱۲۳}$^{*}$ & \lr{۶۷٫۷} & \lr{+۰٫۴۴} & \lr{۰٫۱۴۶}$^{*}$ & \lr{۶۵٫۱} & \lr{+۰٫۱۶} & \lr{۰٫۰۹۱}$^{*}$ & \lr{۵۸٫۹} & \lr{+۰٫۷۶} & \lr{۰٫۱۳۰}$^{*}$ \\
\midrule
\textbf{میانگین} & \lr{۷۵٫۵} & \lr{+۰٫۰۸} & \lr{۰٫۰۸۳} & \lr{۷۰٫۹} & \lr{+۰٫۲۴} & \lr{۰٫۰۹۲} & \lr{۶۶٫۲} & \lr{+۰٫۵۵} & \lr{۰٫۰۸۹} & \lr{۵۵٫۵} & \lr{+۰٫۴۰} & \lr{۰٫۰۹۵} \\
\bottomrule
\end{tabular}
}
\end{table}


\paragraph{یک احتیاط مهم.} رشد $r$ پیش از حرکت با فرضیه‌ی ما می‌خواند، ولی توضیح دیگری هم دارد. در افق بلند، بیشتر خطای \lr{K} پیش از حرکت به یک پرسش برمی‌گردد: «حرکت کی شروع می‌شود؟» آزمایش ۱ نشان داد \lr{E} زمان‌بندی تکلیف را خوب می‌داند و حتی پاسخ به روشن شدن \lr{LED} برای بخش بزرگی از دقتش کافی است. \lr{LED} پیش از حرکت روشن می‌شود. پس \lr{EEG} ممکن است زمان شروع حرکت را از پاسخ بینایی به \lr{LED} یا حرکت چشم به سمت آن بداند، نه از فرمان قشر حرکتی. دو نشانه این توضیح را تقویت می‌کند. اول، اطلاعات فقط در \lr{tokens} آموزش‌دیده حفظ می‌شود و در \lr{band power} خام با افق کم می‌شود. اگر اطلاعات از \lr{ERD} قشر حرکتی بود، \lr{band power} بتا هم باید دست‌کم ثابت می‌ماند. دوم، آزمایش ۲ نشان داد \lr{E} بخشی از اطلاعاتش را از الکترودهای کنار چشم و پشت گردن می‌خواند.

% ---------------------------------------------------------------------------------------------
\FloatBarrier
\subsection{جمع‌بندی سه آزمایش}

\begin{enumerate}
\item \textbf{زمان‌بندی:} دقت \lr{E} در هر ۱۲ نفر از سقف یک مدل زمان‌بندی تکلیف پایین‌تر است و ترتیب عضله‌ها در \lr{E} همان ترتیب مدل زمان‌بندی است. دانستن زمان روشن شدن \lr{LED} به‌تنهایی تا حدود ۹۰ درصد دقت \lr{E} را می‌دهد ($r$ \lr{۰٫۵۸} در برابر \lr{۰٫۶۴}).
\item \textbf{نقشه‌ها:} \lr{ERD} الگوی کلاسیک و تمیزی در هر ۱۲ نفر دارد و سه الکترود اولش (\lr{C3}، \lr{CP1} و \lr{C4}) همان سه الکترود اول \lr{E} است. همبستگی نقشه‌ی \lr{E} با \lr{ERD} بتا و با \lr{CMC} در سطح گروه معنادار ولی متوسط است ($\rho$ حدود \lr{۰٫۲}). \lr{E} از الکترودهای کنار چشم و پشت گردن هم می‌خواند.
\item \textbf{افق:} مدل علّی برای \lr{K} بدون هزینه کار می‌کند. با بزرگ شدن افق، $F-K$ تا حدود نیم واحد \lr{VAF} بالا می‌رود (معنادار در ۵۰۰ میلی‌ثانیه)، و همبستگی پیش‌بینی \lr{EEG} با خطای \lr{K} پیش از حرکت از \lr{۰٫۰۶} به \lr{۰٫۱۹} می‌رسد (رشد معنادار). اما فقط در مرحله‌ی رساندن دست به کاهش واقعی خطا تبدیل می‌شود (تا ۵ درصد، که نیمی از آن به زمان بستگی ندارد).
\end{enumerate}

جمع‌بندی کلی این آزمایش‌ها همراه با فصل‌های قبل در فصل بعد آمده است.
